\subsection[OBSERVATION \& INTELLIGENCE]{\mbox{OBSERVATION \& INTELLIGENCE}}

The Confederates had a rather substantial advantage over the Union during operations in the Shenandoah region due to the fact that a majority of the inhabitants were sympathetic to the Confederacy (a large proportion of the Confederate forces in the area were recruited there). These people could be counted on to pass information to Confederate scouts, cavalry patrols, and headquarters.

As a result, the Confederate army usually had a very good idea of where the Union forces were located, and how many of them there were. The Union army often had great difficulty in gathering intelligence information in this part of the country, usually being forced to rely on the direct observations of their cavalry. This led to the Union armies often not knowing exactly where the Confederates were, or their strength.

\begin{enumerate}
  \item The Confederate Player places the SCREEN around his Stacking Boxes in such a way that the Union Player cannot see the Confederate Stacking Boxes. This allows the Confederate Player to place his units in the Stacking Boxes, and replace them on the Mapboard with Stacking Counters; the Union Player cannot tell what units each Stacking Counter represents.

  \item At any one time, the Confederate Player may have up to six "fake" Stacking Counters on the Map. These "fake" Counters have only Blank units in their appropriate Stacking Boxes. The Confederate Player should try to move thse "fakes" just as they would regular units to prevent the Union Player from guessing that they are "fakes". This will make it difficult for the Union Player to discover exactly where the Confederate units are located.

  \item At all times the Union units are out in the open, and can be inspected by the Confederate Player, as in the Basic Game.

  \item Union units can "observe" the Confederate ones only when two or less hexes from them. When two hexes from a Confederate Counter, the Union Player must be told if it is a real unit, or if it is a "fake". If "fake", the Counter should be immediately removed from the Map. If the Counter contains actual units, the Confederate Player must tell him the types of units in the hex, but now many of them there are. When adjacent to a Confederate Counter, the Union Player determines its exact composition.

  \item New "fake" counters can be placed on the Map, and moved, during any part of the Confederate Player's Phase of the Turn.

  \item When entering new units onto the Mapboard, the Confederate Player may indicate when they are due by placing a Stacking Counter in the appropriate box of the TURN RECORD CHART. If the Confederate Player decided not to bring the units on, the counter could be a "fake". At any rate, it keeps the Union Player guessing. When doing this, the Confederate Player should write down on a piece of paper if the unit is a real one, or a "fake" one, and place the paper in a prominent location. The Union Player may read the note at the end of the game, and adjust the Victory Point totals accordingly.

  \item Units cannot observe through a ridge hexside.
\end{enumerate}